\documentclass[a4paper, serif, 10pt]{moderncv}

%% moderncv
% color and style
\moderncvcolor{black}
\moderncvstyle{banking}

%% page layout/margins
\usepackage[top=2.5cm, bottom=2.5cm, left=1.5cm, right=1.5cm]{geometry}

\nopagenumbers{}

%% Babel config for English language
\usepackage[english]{babel}

%% fontspec
\usepackage{fontspec}

\IfFontExistsTF{Linux Libertine}{
  \setmainfont[Ligatures={TeX, Common}, Numbers=OldStyle]{Linux Libertine}
}{}
\IfFontExistsTF{Linux Libertine O}{
  \setmainfont[Ligatures={TeX, Common}, Numbers=OldStyle]{Linux Libertine O}
}{}

%% CTeX
% CJK support, used to show CJK characters in the resume
%
% - heading=false: disable ctex heading
% - punct=kaiming: use kaiming punctuations styles for CJK
% - scheme=plain: use plain scheme, do not override \normalsize font size
% - space=auto: space settings for CJK characters
%
% CJK fonts are resolved per locale language, e.g. Noto Serif/Sans CJK JP
% for Japanese resumes, with fallback to the next region if not installed
%
% ref:
% - http://ctan.mirrorcatalogs.com/language/chinese/ctex/ctex.pdf
\usepackage[UTF8, heading=false, punct=kaiming, scheme=plain, space=auto]{ctex}

\IfFontExistsTF{Noto Serif CJK SC}{
  \setCJKmainfont{Noto Serif CJK SC}
}{}
\IfFontExistsTF{Noto Sans CJK SC}{
  \setCJKsansfont{Noto Sans CJK SC}
}{}

%% line spacing
\usepackage{setspace}
\setstretch{1.125}

%% URL styling - use normal text instead of monospace
\urlstyle{same}

%% Patch \httplink and \httpslink to handle full URLs
% The original moderncv macros always prepend http:// or https:// to URLs.
% This causes issues when the URL already has a protocol (e.g., https://example.com)
% resulting in malformed URLs like https://https://example.com.
% This patch checks if the URL already contains "://" and uses it directly if so.
% Note: We use \str_set:Nx (x-expansion) to expand macros like \@homepage first.
\ExplSyntaxOn
\str_new:N \g_yamlresume_url_str

\RenewDocumentCommand{\httpslink}{O{}m}{
  \str_set:Nx \g_yamlresume_url_str {#2}
  \str_if_in:NnTF \g_yamlresume_url_str {://}
    { \url{#2} }
    { \url{https://#2} }
}

\RenewDocumentCommand{\httplink}{O{}m}{
  \str_set:Nx \g_yamlresume_url_str {#2}
  \str_if_in:NnTF \g_yamlresume_url_str {://}
    { \url{#2} }
    { \url{http://#2} }
}
\ExplSyntaxOff

\name{Morgan Casey}{}
\title{DevOps Engineer | Cloud Infrastructure \& Automation}
\phone[mobile]{(512) 555-0147}
\email{morgan.casey@example.com}
\homepage{https://morgancasey.dev}

\address{2400 Nueces Street, Austin, Texas, United States, 78705}{}{}

\extrainfo{{\small \faGithub}\ \href{https://github.com/morgancasey}{@morgancasey} {} {} {} • {} {} {} 
{\small \faLinkedin}\ \href{https://www.linkedin.com/in/morgancasey}{@morgancasey}}

\begin{document}

\maketitle

\section{Basics}

\cvline{}{DevOps engineer with 6+ years of experience designing, automating, and operating cloud-native infrastructure on AWS and Azure. Skilled in building CI/CD pipelines, implementing infrastructure as code, and orchestrating containerized workloads with Kubernetes. Proven record of cutting deployment times, improving system reliability, and building strong collaboration between development and operations teams.
%
\begin{itemize}
\item Reduced average deployment time from 45 minutes to 8 minutes by migrating legacy Jenkins pipelines to GitLab CI/CD.
\item Managed 200+ microservices across multiple Kubernetes clusters using Helm, Argo CD, and Terraform.
\item Implemented an observability stack with Prometheus, Grafana, and the ELK Stack, cutting mean time to detection by 40\%.
\end{itemize}}
\section{Education}

\cventry{Aug 2014–May 2018}
        {Bachelor, Computer Science, Score: 3.7}
        {University of Texas at Austin}
        {\url{https://www.utexas.edu/}}
        {}
        {\begin{itemize}
\item Focused on distributed systems, networking, and cloud computing
\item Completed a capstone project on container orchestration and cluster scheduling
\item Served as a teaching assistant for the Operating Systems course for two semesters
\end{itemize}
\textbf{Courses}: Operating Systems, Computer Networks, Distributed Systems, Database Systems, Cloud Computing, Algorithms and Data Structures}
\section{Work}

\cventry{Jan 2022–Present}
        {Senior DevOps Engineer}
        {CloudForge Systems}
        {\url{https://www.cloudforgesystems.com}}
        {}
        {\begin{itemize}
\item Lead a team of 4 DevOps engineers managing multi-region AWS infrastructure for a SaaS platform serving 500K+ users
\item Architected GitOps workflows with Argo CD and Terraform, enabling self-service deployments for 12 product teams
\item Reduced cloud spend by 28\% through right-sizing, reserved instances, and automated scheduling of non-production environments
\item Designed a disaster recovery strategy with cross-region backups and automated failover, achieving an RTO of 15 minutes and an RPO of 5 minutes
\item Mentored junior engineers and established on-call best practices, decreasing escalations by 35\%
\end{itemize}
\textbf{Keywords}: AWS, Terraform, Kubernetes, Argo CD, GitOps, Cost Optimization}

\cventry{Jun 2019–Dec 2021}
        {DevOps Engineer}
        {NimbusStack}
        {\url{https://www.nimbusstack.io}}
        {}
        {\begin{itemize}
\item Built and maintained CI/CD pipelines in GitLab CI for 80+ microservices, integrating automated testing and security scans
\item Migrated a monolithic application to Kubernetes, improving scalability and reducing infrastructure costs by 20\%
\item Implemented infrastructure as code with Terraform and Ansible across 300+ EC2 instances and RDS databases
\item Set up centralized logging and monitoring with the ELK Stack and Prometheus, enabling proactive alerting
\item Partnered with developers to containerize applications and optimize Docker images for faster cold starts
\end{itemize}
\textbf{Keywords}: GitLab CI, Kubernetes, Terraform, Ansible, Prometheus, Docker}

\cventry{Jul 2018–May 2019}
        {Cloud Operations Engineer}
        {DataHarbor}
        {\url{https://www.dataharbor.com}}
        {}
        {\begin{itemize}
\item Managed AWS EC2, S3, and RDS resources for an internal data analytics platform
\item Automated routine operational tasks with Python and Bash, saving roughly 15 hours of manual effort per week
\item Assisted in implementing backup and recovery procedures for critical production databases
\item Participated in the on-call rotation and incident response for production systems
\end{itemize}
\textbf{Keywords}: AWS, Python, Bash, Linux, Backup}
\section{Languages}

\cvline{English}{Native or Bilingual Proficiency \hfill \textbf{Keywords}: TOEFL 115}
\cvline{Spanish}{Limited Working Proficiency \hfill \textbf{Keywords}: DELE B2}
\section{Skills}

\cvline{Cloud Platforms}{Expert \hfill \textbf{Keywords}: AWS, Azure, GCP, EC2, S3, Lambda}
\cvline{CI/CD}{Advanced \hfill \textbf{Keywords}: GitLab CI, Jenkins, GitHub Actions, Argo CD, CircleCI}
\cvline{Containers \& Orchestration}{Expert \hfill \textbf{Keywords}: Kubernetes, Docker, Helm, Istio, OpenShift}
\cvline{Infrastructure as Code}{Advanced \hfill \textbf{Keywords}: Terraform, Ansible, CloudFormation, Pulumi}
\cvline{Monitoring \& Logging}{Advanced \hfill \textbf{Keywords}: Prometheus, Grafana, ELK Stack, Datadog, CloudWatch}
\cvline{Scripting \& Automation}{Expert \hfill \textbf{Keywords}: Python, Bash, Go, PowerShell}
\cvline{Security \& Compliance}{Intermediate \hfill \textbf{Keywords}: IAM, Vault, SOC 2, CIS Benchmarks, Trivy}
\section{Awards}

\cventry{Nov 2023}
        {CloudForge Systems}
        {Cloud Innovator Award}
        {}
        {}
        {Recognized for designing and implementing a cost optimization program that saved more than \$500K annually.}
\section{Certificates}

\cventry{Mar 2022}
        {Amazon Web Services}
        {AWS Certified DevOps Engineer - Professional}
        {\url{https://aws.amazon.com/certification/certified-devops-engineer-professional/}}
        {}
        {}

\cventry{Aug 2021}
        {Cloud Native Computing Foundation}
        {Certified Kubernetes Administrator (CKA)}
        {\url{https://www.cncf.io/certification/cka/}}
        {}
        {}

\cventry{Jan 2021}
        {HashiCorp}
        {HashiCorp Certified: Terraform Associate}
        {\url{https://www.hashicorp.com/certification/terraform-associate}}
        {}
        {}
\section{Publications}

\cventry{Sep 2023}
        {Implementing GitOps for Multi-Cluster Kubernetes}
        {DevOps Today}
        {\url{https://www.devopstoday.com/gitops-multi-cluster-kubernetes}}
        {}
        {\begin{itemize}
\item Discusses patterns for managing multiple Kubernetes clusters with Argo CD and Terraform
\item Covers security, drift detection, and progressive delivery strategies
\end{itemize}}
\section{Projects}

\cventry{Feb 2022–Nov 2023}
        {Open-source Kubernetes monitoring and alerting toolkit}
        {KubeWatch}
        {\url{https://github.com/morgancasey/kubewatch}}
        {}
        {\begin{itemize}
\item Built a lightweight operator that watches cluster events and routes alerts to Slack and PagerDuty
\item Integrated with Prometheus Alertmanager for advanced routing and silencing
\item Reached 500+ stars on GitHub and is used by more than 30 organizations
\end{itemize}
\textbf{Keywords}: Kubernetes, Go, Prometheus, Operator}

\cventry{Mar 2021–Aug 2022}
        {Reusable Terraform modules for secure AWS multi-account setups}
        {Terraform AWS Landing Zone}
        {\url{https://github.com/morgancasey/terraform-aws-landing-zone}}
        {}
        {\begin{itemize}
\item Published a set of Terraform modules implementing AWS Control Tower best practices
\item Automated account provisioning, guardrails, and centralized logging
\item Reduced new account setup time from days to under an hour
\end{itemize}
\textbf{Keywords}: Terraform, AWS, Landing Zone, Security}
\section{Interests}

\cvline{Open Source}{CNCF, Kubernetes, Terraform}
\cvline{Outdoor Activities}{Hiking, Running, Cycling}
\section{Volunteer}

\cventry{Mar 2020–Present}
        {Workshop Organizer}
        {Austin DevOps Community}
        {\url{https://www.meetup.com/austin-devops/}}
        {}
        {\begin{itemize}
\item Organize monthly workshops on cloud-native technologies and DevOps practices
\item Coordinate with speakers and sponsors to deliver free technical education
\item Mentor newcomers transitioning into DevOps roles
\end{itemize}}
    
\end{document}